\section{Discussion}

\label{sec:discussion}

\subsection{Vendor Lock-in Mitigation}

Hanzo Platform uses standard Kubernetes primitives and OCI containers, enabling applications to be extracted and deployed on any Kubernetes cluster. The \texttt{hanzo eject} command generates standard Kubernetes manifests, Dockerfiles, and Helm charts for migration.

\subsection{Limitations}

\begin{enumerate}[leftmargin=1.4em]
    \item \textbf{GPU availability}: GPU capacity is finite; burst scaling may be delayed during high-demand periods.
    \item \textbf{Cold starts}: GPU inference cold starts (30--120s for model loading) remain a challenge despite pre-warming.
    \item \textbf{Regional coverage}: Currently deployed in 3 regions (US East, US West, EU West); Asia-Pacific coverage is planned.
    \item \textbf{Custom hardware}: Only NVIDIA GPUs are currently supported; AMD and custom accelerators are on the roadmap.
\end{enumerate}

\subsection{Future Work}

\begin{itemize}[leftmargin=1.1em]
    \item \textbf{Edge inference}: Deploy small models at CDN edge nodes for ultra-low-latency inference.
    \item \textbf{Multi-cloud}: Support AWS, GCP, and Azure as compute backends alongside DigitalOcean.
    \item \textbf{Fine-tuning as a service}: Managed fine-tuning with automatic deployment of resulting models.
    \item \textbf{AI-assisted operations}: Use LLMs to diagnose deployment failures and suggest fixes.
\end{itemize}

